%% ==========================================================================
%%  FINAL REPORT --- SECTION 6, EVALUATION
%%  Every number below is recomputable from the newest-sha rows of the named
%%  bench/results/*.jsonl file, the same filter make_figures.py applies.
%%  Figures F3 F4 F5 F7 F8 F9 F12 live here, each exactly once.
%% ==========================================================================
\section{Evaluation}
\label{sec:eval}

This section asks what private training and private delivery each cost, what
they cost together, where the cost goes, and what it does to recommendation
quality. It reports two results that go against the system in the body rather
than in a footnote.

\subsection{Setup}
\label{sec:eval:setup}

All measurements use MovieLens-100K's \texttt{u1} split (943 users, 1{,}682
items, 80{,}000 training ratings) at $\dime = 16$, $\iters = 10$ and $\fracbits =
20$, with $\bits = 64$ and the revealing normaliser unless stated. Delivery
fetches $\topk = 10$ records per session, and quality is nDCG@20. The machine is
a Windows~11 laptop, and the code is built by MSYS2's g++.

\paragraph{Network profiles.} We use four: \texttt{local}, \texttt{lan} (1\,ms
round trip, 1\,Gbit/s), \texttt{wan\_a} (30\,ms, 100\,Mbit/s) and \texttt{wan\_b}
(100\,ms, 10\,Mbit/s). The development machine has neither Docker nor kernel
network emulation, so the profiles are applied at the channel level. Delivery
runs over real sockets wrapped in a channel that delays each frame. Training runs
its three parties in one process, so its time on a profile is computed from the
exactly counted rounds and bytes as rounds times round-trip time plus bytes over
bandwidth. We call this channel-level emulation throughout and never kernel
emulation.

\paragraph{What to trust.} Round and byte counts are exact and deterministic.
Absolute timings on this laptop vary by up to a factor of four between
back-to-back runs, while repetitions within one run agree closely, which is the
pattern of thermal throttling. So we quote ratios taken within one run, and where
a number derived from absolute times moves between runs we give the range we
observed. Every figure is drawn by \texttt{mingw32-make figures} from the newest
commit's rows in the named data file, and Appendix~\ref{app:reproduce} maps each
one to the command that produces it.

\subsection{Recommendation quality}
\label{sec:eval:quality}

Table~\ref{tab:quality} compares private training with the cleartext oracle, which
runs the same power iteration in floating point. The revealing path at
$\bits = 64$ stays within 0.007 of the oracle at every $\iters$. Fixed-point
arithmetic and a real truncation protocol cost essentially nothing in quality.

\begin{table}[t]
\centering
\small
\caption{nDCG@20 against the number of power iterations $\iters$, at $\dime = 16$.
The $\bits = 128$ revealing column is the control that isolates the normaliser:
it shares the spec-faithful path's ring and truncation schedule and differs only
in revealing the norm. The last column is the no-leak path's loss against that
control. Data: \texttt{bench/results/train\_quality.jsonl}.}
\label{tab:quality}
\begin{tabular}{@{}rccccc@{}}
\toprule
$\iters$ & reveal, $\bits = 64$ & reveal, $\bits = 128$ & \textbf{no leak},
  $\bits = 128$ & oracle & no-leak cost \\
\midrule
1  & 0.3729 & 0.3727 & \textbf{0.1651} & 0.3664 & 0.2076 \\
2  & 0.4251 & 0.4248 & 0.4230 & 0.4220 & 0.0019 \\
5  & 0.4485 & 0.4490 & 0.4448 & 0.4455 & 0.0041 \\
10 & 0.4511 & 0.4513 & 0.4493 & 0.4536 & 0.0020 \\
25 & 0.4532 & 0.4531 & 0.4499 & 0.4545 & 0.0031 \\
50 & 0.4531 & 0.4532 & 0.4502 & 0.4525 & 0.0030 \\
\bottomrule
\end{tabular}
\end{table}

\paragraph{Revealing nothing costs little quality.} For every $\iters \ge 2$ the
spec-faithful normaliser loses between 0.0018 and 0.0041 of nDCG@20 against
either revealing column. The $\bits = 128$ revealing column is what makes that
claim defensible. Without it, the comparison would confound the normaliser with
the ring width and the deferred truncation schedule that $\bits = 128$ allows.
The two revealing columns differ by at most 0.0005, so the normaliser accounts
for the whole gap. What revealing nothing costs instead is rounds: 11{,}519
against 2{,}191 at $\iters = 10$ and $\bits = 128$, a factor of 5.3.

\paragraph{The exception at $\iters = 1$.} With a single power iteration the
no-leak path scores 0.1651 against 0.3729. The first normalisation acts on the
raw random start vector, and a single iteration leaves nothing downstream to
recover from whatever error it introduces. We have not isolated the mechanism
further. From $\iters = 2$ the gap closes. Nobody would choose $\iters = 1$, but
the cliff is sharp and we report it.

\paragraph{Quality is flat in $\iters$.} The oracle's own sweep settles how many
iterations are worth paying for (Figure~\ref{fig:ell}). Across a fortyfold
increase in $\iters$ the subspace angle to the true singular subspace falls by
five orders of magnitude, from 0.97 radians to $6.7 \times 10^{-6}$, while
nDCG@20 stays between 0.4523 and 0.4545. Ranking needs a subspace good enough to
order the scores stably, not the singular vectors themselves. The slow
convergence has a visible cause: the sixteenth and seventeenth squared singular
values differ by only 3\%. Since rounds grow linearly in $\iters$
(Section~\ref{sec:eval:sweep}), this flatness is what makes $\iters = 10$ the
right operating point.

\begin{figure}[t]
\centering
\includegraphics[width=0.62\linewidth]{fig_ell}
\caption{The cleartext oracle at $\iters$ from 10 to 400. nDCG@20 spans 0.002
while the subspace angle spans five orders of magnitude. Data:
\texttt{bench/results/oracle\_ell\_sweep.jsonl}.}
\label{fig:ell}
\end{figure}

\subsection{The cost of each half, isolated}
\label{sec:eval:compose}

The central comparison runs five arms in one process, interleaved repetition by
repetition so that thermal drift cannot line up with arm identity
(Table~\ref{tab:compose} and Figure~\ref{fig:compose}). B1 trains in the clear
and fetches in the clear. B2 swaps in DPF-PIR. B3 swaps in private training
instead. FULL is \sys{}. B5 trains in the clear and downloads the whole catalogue.
Training is amortised over the 943 users one training run serves.

\begin{table}[t]
\centering
\small
\caption{Per-session cost of each arm, with training amortised over 943 users.
Bytes and rounds are exact. Times are channel-level emulation. Data:
\texttt{bench/results/bench\_compose.jsonl}.}
\label{tab:compose}
\begin{tabular}{@{}llrrrrc@{}}
\toprule
Arm & training, delivery & bytes & rounds & \texttt{local} (ms) &
  \texttt{wan\_a} (ms) & nDCG@20 \\
\midrule
B1   & clear, clear         & 2{,}600   & 10 & 1.19 & 301.4 & 0.454 \\
B2   & clear, DPF-PIR       & 9{,}660   & 10 & 3.11 & 303.9 & 0.454 \\
B3   & private, clear       & 71{,}264  & 13 & 2.02 & 397.7 & 0.451 \\
FULL & private, DPF-PIR     & 78{,}324  & 13 & 3.96 & 400.2 & 0.451 \\
B5   & clear, full download & 430{,}592 & 1  & 1.16 & 65.6  & 0.454 \\
\bottomrule
\end{tabular}
\end{table}

\paragraph{Composing costs exactly the sum of the halves.} $\mathrm{FULL} =
\mathrm{B2} + \mathrm{B3} - \mathrm{B1}$ holds to the byte. Over a whole training
run plus one session it is 64{,}760{,}124 bytes, and the benchmark asserts the
identity before it writes a single row, so a regression that broke additivity
could not produce a figure. On \texttt{wan\_a}, private delivery adds 2.5\,ms per
session and private training adds 96\,ms, and quality moves from 0.454 to 0.451.

\paragraph{The result that goes against us.} On \texttt{wan\_a}, B5 is 4.6 times
\emph{faster} than B2, although it sends 45 times as many bytes. Delivery as
built costs one round trip per recommendation, ten in all, while B5 pays one. At
30\,ms per trip, round trips decide the outcome and bytes barely matter.
Sending all ten key pairs in one round trip reverses it: the model puts batched
B2 at 33.9\,ms against B5's 65.6\,ms. We did not implement batching, and the
benchmark marks those rows as not implemented. The honest summary is that
private delivery as built is cheap in bytes and expensive in round trips, and the
fix is known and unbuilt.

\begin{figure}[t]
\centering
\includegraphics[width=\linewidth]{fig_compose}
\caption{Left: what each arm adds over B1, per session, on \texttt{local} and
\texttt{wan\_a}. Right: the \texttt{wan\_a} session time as built, against the
modelled time with the $\topk$ queries batched into one round trip, which is not
implemented. Data: \texttt{bench/results/bench\_compose.jsonl}.}
\label{fig:compose}
\end{figure}

\subsection{The private read against its alternatives}
\label{sec:eval:pir}

Figure~\ref{fig:pircost} compares one private read with a cleartext lookup (B1)
and with downloading the whole catalogue (B5), the first alternative a reader
proposes because 430\,KB does not sound large. DPF-PIR sends 483 bytes to each
server, so it moves 446 times fewer bytes than B5 counting both servers. It is
also the slowest of the three in server CPU, about five times slower than B5's
local copy in both committed runs (5.3 and 5.6), and four orders of magnitude
slower than a cleartext lookup. Its advantage is entirely in communication.

Setting the extra server time against the saved bytes gives the link speed below
which the private read wins. That crossover is derived from absolute times, so it
moves with the machine's thermal state: 8.0\,Gbit/s in one committed run and
22.9\,Gbit/s in the run the figure shows. Either way it is far above any
wide-area link, and we would rather state the crossover than imply there is none.

\begin{figure}[t]
\centering
\includegraphics[width=0.9\linewidth]{fig_pir_cost}
\caption{One read at the MovieLens-100K operating point: 1{,}682 records of
256\,B in an 11-bit domain. Bytes are per server. Both panels must be read
together, since the left alone would be advocacy. Data:
\texttt{bench/results/bench\_baseline.jsonl}.}
\label{fig:pircost}
\end{figure}

\subsection{Scaling, and Theorem~4.2 measured}
\label{sec:eval:sweep}

The sweep varies one parameter at a time around the baseline, five repetitions
each, on all four profiles (Figure~\ref{fig:sweep}). It is not the Cartesian
product, so interactions between parameters are not measured.

\paragraph{Rounds.} Rounds are linear in $\dime$ (1{,}310, 2{,}686 and 5{,}438
at $\dime = 8, 16, 32$) and in $\iters$ (436 at $\iters = 1$, 6{,}436 at
$\iters = 25$), and flat in $\users$: 2{,}686 at 235, 471 and 943 users.

\paragraph{Bytes.} This is \nudgeplain{}'s Theorem~4.2 measured rather than
quoted. Quadrupling the input matrix, from 235 to 943 users, moves traffic from
62.03 to 64.75\,MB, an increase of 4.4\%. The increase is exactly linear at
3{,}840 bytes per added user, which is the one term in the cost that depends on
$\users$: each of the $\dime\iters = 160$ iterations computes
$\ratings\mathbf{v}$, an $\users$-vector, and resharing it costs three parties
one 8-byte element per entry, so $160 \times 3 \times 8 = 3{,}840$ bytes per
user. Communication tracks the largest intermediate vector, not the input matrix.

\paragraph{Training is round-bound.} On \texttt{wan\_a} one training run takes
85.8\,s, of which 80.6\,s is round-trip latency and 5.2\,s is transfer. Anything
that removes rounds matters far more than anything that removes bytes, which is
why the next subsection counts rounds.

\begin{figure}[t]
\centering
\includegraphics[width=\linewidth]{fig_sweep}
\caption{Left: rounds against $\dime$, $\iters$ and $\users$. Right: traffic
against the size of the input matrix. The matrix grows fourfold and traffic by
4.4\%. Data: \texttt{bench/results/bench\_sweep.jsonl}.}
\label{fig:sweep}
\end{figure}

\subsection{Where the rounds go}
\label{sec:eval:stages}

Table~\ref{tab:stages} attributes every round and byte of one training run to
the stage that spent it, in the three configurations. The attribution is a closed
form derived from each protocol's structure, then validated against a real run:
the residual is zero rounds and zero bytes in all three configurations, and the
benchmark exits with an error if it is not. So these are shares of a complete
budget, not a sample. A zero residual alone would not catch rounds attributed to
the wrong stage, so a second, independent count checks that too. A wrapper around
the normaliser meters every round spent inside it, and the model must predict
that count exactly, including the truncations the normaliser makes internally.
It does in all three configurations: 880 rounds for each revealing path and
10{,}208 for the path that reveals nothing.

\begin{table}[t]
\centering
\small
\caption{Rounds per training run by stage, with each stage's share of rounds and
of bytes. ``Open $\itemembed$'' is the sixteen row openings. Data:
\texttt{bench/results/bench\_stages.jsonl}.}
\label{tab:stages}
\begin{tabular}{@{}lrrrrrr@{}}
\toprule
 & \multicolumn{2}{c}{reveal, $\bits = 64$} & \multicolumn{2}{c}{reveal,
 $\bits = 128$} & \multicolumn{2}{c}{no leak, $\bits = 128$} \\
\cmidrule(lr){2-3}\cmidrule(lr){4-5}\cmidrule(l){6-7}
Stage & rounds & \% rounds / bytes & rounds & \% rounds / bytes & rounds &
  \% rounds / bytes \\
\midrule
matrix--vector & 320 & 11.9 / 15.6 & 320 & 14.6 / 15.6 & 320 & 2.8 / 14.0 \\
truncation & 1{,}998 & \textbf{74.4} / \textbf{83.4} & 1{,}503 & \textbf{68.6} /
  \textbf{83.4} & 8{,}367 & \textbf{72.6} / \textbf{75.1} \\
normalisation & 352 & 13.1 / 0.0 & 352 & 16.1 / 0.0 & 2{,}464 & 21.4 / 9.9 \\
FSS gate & 0 & 0 / 0 & 0 & 0 / 0 & 352 & 3.1 / 0.0 \\
open $\itemembed$ & 16 & 0.6 / 1.0 & 16 & 0.7 / 1.0 & 16 & 0.1 / 0.9 \\
\midrule
total & 2{,}686 & & 2{,}191 & & 11{,}519 & \\
\bottomrule
\end{tabular}
\end{table}

\paragraph{Truncation dominates, in every configuration.} It is 68.6 to 74.4\%
of rounds and 75.1 to 83.4\% of bytes. The FSS comparison gate, in the one
configuration that has it, is 3.1\% of rounds. The shared matrix--vector products
cost 320 rounds, one per product, and between 2.8 and 14.6\% of rounds.
Anyone optimising this system should start with truncation, whose three rounds
recur after every rescaling, every orthogonalisation and every Newton step. The
gate, which looks like the exotic part, is not the problem.

\paragraph{CPU ranks the stages differently.} Figure~\ref{fig:stages} also times
each operation on its own. The shared matrix--vector product over the full rating
matrix is the most expensive operation in CPU by a wide margin, while being cheap
in rounds. On a wide-area link the round count governs, and locally the CPU time
does, so which stage matters depends on where the servers are.

\begin{figure}[t]
\centering
\includegraphics[width=\linewidth]{fig_stages}
\caption{Left: each stage's share of the rounds in one training run, per
configuration. Right: median CPU time per call, each operation timed on its own.
Data: \texttt{bench/results/bench\_stages.jsonl}.}
\label{fig:stages}
\end{figure}

\subsection{How wide the ring must be}
\label{sec:eval:ring}

Whether a 64-bit ring suffices at MovieLens scale, or \nudgeplain{}'s 128 bits are
needed, was an open question in our design. It has two answers, for two different
reasons, and the tighter one wins.

\paragraph{Arithmetic headroom.} The truncation schedule is derived from public
parameters, and deriving it at larger scales answers the question without running
them. At $\fracbits = 20$ and $\bits = 64$, the deferred schedule does not fit at
any of the three scales we evaluated, MovieLens-100K, MovieLens-1M and a
Netflix-sized matrix, so training truncates immediately, with 21, 18 and 11 bits
to spare. At $\fracbits = 24$ and Netflix size there is no safe schedule at all.
At $\bits = 128$ the deferred schedule fits everywhere with at least 55 spare
bits, and at MovieLens-100K it needs fewer rounds (2{,}191 against 2{,}686) for
twice the bytes (129.3\,MB against 64.75\,MB).

\paragraph{The mask.} The spec-faithful normaliser's gate needs an 81-bit domain
to hide its input (Section~\ref{sec:impl:norm}). So a 64-bit ring is enough
arithmetically at MovieLens-1M, and not enough at any scale for a normaliser that
reveals nothing. That second constraint is not arithmetic, and it is the one that
decides the ring width.

\subsection{On the wire}
\label{sec:eval:wire}

An argument that message sizes do not depend on the fetched index, and a
measurement that they do not, are different claims. We recorded every frame at
the channel boundary while a probe issued 1{,}800 queries across three record
indices, 0, 777 and 1234, generating fresh keys for every query. Every query frame
is 228 bytes and every reply is 265 bytes, whatever the index.

We then trained a logistic-regression classifier on the raw frame bytes to tell
each pair of indices apart and scored it on 480 held-out queries per pair
(Table~\ref{tab:distinguisher} and Figure~\ref{fig:distinguisher}). Chance lies
inside every 95\% interval. The same pipeline run on deliberately shuffled labels
agrees with the real runs, which shows the harness measures what it claims to.

\begin{table}[t]
\centering
\small
\caption{Telling fetched indices apart from the wire bytes alone, with 480
held-out queries per pair and the same pipeline run on shuffled labels as a
control. Data: \texttt{bench/results/distinguisher\_result.jsonl}.}
\label{tab:distinguisher}
\begin{tabular}{@{}lcc@{}}
\toprule
Adversary & accuracy & 95\% interval \\
\midrule
index 0 against 1234   & 0.490 & 0.445 to 0.534 \\
index 0 against 777    & 0.494 & 0.449 to 0.538 \\
index 1234 against 777 & 0.500 & 0.455 to 0.545 \\
\midrule
shuffled labels        & 0.483 & 0.439 to 0.528 \\
\bottomrule
\end{tabular}
\end{table}

\begin{figure}[t]
\centering
\includegraphics[width=0.62\linewidth]{fig_distinguisher}
\caption{The distinguisher at chance. What matters is that 0.5 lies inside
every interval, not that the points sit near it. Data:
\texttt{bench/results/distinguisher\_result.jsonl}.}
\label{fig:distinguisher}
\end{figure}

This is not a proof. It says that one adversary, with one feature representation,
did not beat chance on this data. What it rules out is the failure that actually
occurs in practice, a length or byte pattern that correlates with the query.

\subsection{Server-side timing}
\label{sec:eval:timing}

Frame sizes are fixed, but a server whose running time depended on the index
would leak it to anyone who can time the reply. We timed the PIR answer routine
200 times at each of eight indices spread across the domain, interleaved rather
than in blocks so that thermal drift could not line up with the index
(Figure~\ref{fig:timing}). The mean call takes 87\,$\mu$s. The spread between
the eight per-index means is 1.7\,$\mu$s, smaller than the median standard
deviation of 6.0\,$\mu$s within a single index, and a permutation test on
shuffled index labels puts the observed spread at $p = 0.60$, in the middle of
the null distribution.

This is a bound, not a proof of constant time, because no finite sample shows a
difference is exactly zero. What it shows is that any index-dependent signal is
below 2.0\% of the mean call at this sample size. That is consistent with the
construction, which expands every node and takes an unconditional inner product
over every record. Timing on the wire itself is dominated by the operating
system's scheduler and network stack, and we did not analyse it.

\begin{figure}[t]
\centering
\includegraphics[width=\linewidth]{fig_timing}
\caption{Left: server time per call at eight record indices. Right: the spread
between index means, against its distribution under shuffled labels. Data:
\texttt{bench/results/bench\_timing.jsonl}.}
\label{fig:timing}
\end{figure}
